\documentclass[11pt]{article}
\usepackage[margin=1.1in]{geometry}
\usepackage{amsmath,amssymb,amsthm}
\usepackage{booktabs}
\usepackage{array}
\usepackage{tabularx}
\usepackage[protrusion=true,expansion=false]{microtype}
\usepackage[table]{xcolor}
\usepackage{enumitem}
\usepackage{titlesec}
\usepackage{tikz}
\usepackage{float}
\usetikzlibrary{arrows.meta,positioning,calc}
\usepackage{hyperref}

\hypersetup{colorlinks=true, linkcolor=blue!60!black, urlcolor=blue!60!black, citecolor=blue!60!black}
\usepackage[skip=0.6\baselineskip, indent=0pt]{parskip}

\titleformat{\section}{\large\bfseries}{\thesection.}{0.7em}{}
\titleformat{\subsection}{\normalsize\bfseries}{\thesubsection.}{0.6em}{}

\newcommand{\vmax}{V_{\max}}
\newcommand{\Tc}{T_{\mathrm{close}}}
\newcommand{\tx}[1]{\texttt{#1}}
\newcolumntype{L}{>{\raggedright\arraybackslash}X}

\renewcommand{\arraystretch}{1.3}
% float pages: floats from the top, not spread over the page
\makeatletter
\setlength{\@fptop}{0pt}
\setlength{\@fpsep}{1.5\baselineskip}
\setlength{\@fpbot}{0pt plus 1fil}
\makeatother
\rowcolors{1}{}{black!5}

\definecolor{coop}{RGB}{232,245,233}
\definecolor{chainc}{RGB}{255,228,225}
\definecolor{hubc}{RGB}{255,243,214}

\title{\textbf{A ledger for Lichen}\\validity, finality and exit when the hub is the sequencer\\[2pt]\large draft}
\author{waxwing}
\date{Version 0.1, October 2026\\[2pt]\small a companion to \emph{Lichen}, version 0.3~\cite{lichen}}

\begin{document}
\maketitle

\begin{abstract}
Lichen~\cite{lichen} lets a user withdraw a hub's L2 tokens as bitcoin
without trusting the hub, given a proof that her withdrawal is correct
by the L2's own account of its ledger. Like the BitVM bridges, the main
paper leaves that account to the L2. This note describes lichen's own
L2, in which the hub is also the sequencer, so that the party that pays
out controls the ledger it is judged by. That merger forces the
withdrawal statement to enforce what a separate L2 would vouch for: that
every block was valid, that there is one history and that it is final on
Bitcoin, and that a user can leave from a state already anchored even if
the operator censors her or vanishes. We describe the ledger, a chain of
validity proofs, anchoring by a chain of single-use seals with the
contract's work on top, and the exit: Alice's claim from an anchored
state, the hub's counter to it, and the statement's three branches, all
riding on the same BABE gadget as the main paper's. All of it is built.
The cooperative path is unchanged. An exit's dispute is about 15~kvB,
against 9.7 for a refused withdrawal, and the hub's proof of a counter
takes minutes. The worst case falls on Alice:
proving a forged counter false verifies a Groth16 proof inside RISC
Zero, 56.3M cycles, about an hour on a laptop. It needs none of her
secrets and has days to run, so it can be delegated, but it is a real
burden on an ordinary user, and we set out the alternatives. Data
availability remains the main open problem.
\end{abstract}

\tableofcontents

\section{Scope}
\label{sec:scope}

The main paper~\cite{lichen} describes a contract between a user, Alice,
and a hub: a 2-of-2 output holding her deposit, the hub's liquidity up to
a cap $\vmax$, and two reserves. The hub's tokens on an L2 are
collateralised by it; a withdrawal of $b$ is a return of $b$ tokens to
the hub with the contract's memo $c$, folded cooperatively or, if the hub
refuses or is offline, enforced on Bitcoin by a verification gadget,
BABE~\cite{babe} with the amount authenticated by the hub's one-time
signature~\cite{lichenbabe}. What the gadget enforces is a statement
fixed when the contract opens. At its simplest that statement is:
\begin{quote}
\emph{a return of $b$ to contract $c$ exists, by the L2's account of its
ledger.}
\end{quote}
How the L2 makes its account trustworthy is its own business, as it is
for the BitVM bridges~\cite{bitvm2usenix}: an L2 secured by its own
consensus, or a rollup on another chain with its own proofs, supplies
the account and lichen consumes it. The hub is then only a counterparty;
it cannot rewrite a ledger it does not run.

This note describes the L2 that lichen builds~\cite{lichencode}. In it
the hub and the sequencer are one operator, or two that trust each
other. That is a design choice, and its consequence is that the L2's
account is the hub's own word. A hub that profits from
defaults, since whatever isn't withdrawn by $\Tc$ is its, has every
reason to bend that account. So everything a separate L2 would vouch for
has to be enforced inside the withdrawal statement, against the
operator. We build it in three levels:
\begin{itemize}[leftmargin=1.5em]
\item $P_0$, the main paper's: the account is trusted.
\item $P_1$, \emph{validity and finality}: the return is in a state that
  ends a valid history from genesis, that history is the one anchored on
  Bitcoin, and the anchor has the contract's work on top
  (Sections~\ref{sec:validity}--\ref{sec:withdrawal}).
\item $P_2$, \emph{the exit}: Alice can leave from a state already
  anchored, without a return, against an operator that censors her
  return, stops sealing or vanishes; the hub may answer only with a
  counter showing that she spent what she claims
  (Sections~\ref{sec:exit}--\ref{sec:worst}).
\end{itemize}
The gadget does not change between the levels; the statement does, and
from $P_2$ on, so do the first two transactions of the dispute, which
carry the claim and the counter. Section~\ref{sec:open} lists what is
left, data availability first.

Figures marked (lichen) are measured in lichen's proof of concept, on
regtest under mainnet's relay policy or on Mutinynet as stated; the rest
are estimates, marked (est.). Proving times are on one 14-core Apple M4
Pro laptop, on its CPU (the RISC Zero version we use has no Metal
prover), or on a rented 4-vCPU x86 server, as stated.

\section{The ledger}
\label{sec:ledger}

The ledger is one of notes. A note is an owner, a value, a memo and a
random $\rho$; its commitment is a SHA-256 hash of these, and the
commitments form an append-only Merkle tree, so a note's path under a
root stays valid however the ledger grows. A spent note is not removed
but recorded by its \emph{nullifier} in a second tree, a sparse Merkle
tree keyed by nullifier, whose leaf records how the note was spent:
\begin{itemize}[leftmargin=1.5em]
\item \emph{spent}: an ordinary transaction;
\item \emph{returned}$(c, r)$: a transaction paying the hub exactly one
  note, of $r$ with memo $c \ne 0$, the first step of a withdrawal from
  contract $c$;
\item \emph{exited}$(c)$: contract $c$'s exit record
  (Section~\ref{sec:exit}).
\end{itemize}
The tags matter only to the exit, whose counter must tell which spends
may be held against which claim. A state root $\mathit{SR}$ hashes the
height, the commitment tree's root and the nullifier tree's root, and
the sequencer signs it (BIP340).

There are two kinds of note in one tree. \emph{Transparent} notes have a
public owner key and value: the hub's mints (a deposit's tokens) and
returns to the hub are transparent, and a transparent spend is
authorised by the owner's signature. \emph{Shielded} notes publish only
their commitment; the owner's authority is knowledge of a key
$\mathit{nk}$, the nullifier is derived from it so that nobody can link
it to its note, and the sender proves the spend in zero knowledge on her
own machine~\cite{shielded}. The transaction types are a mint, a
transfer, shielding and unshielding, a return, an anchor record
(Section~\ref{sec:finality}) and an exit record (Section~\ref{sec:exit}).

Three proofs, all RISC Zero, are involved (Table~\ref{tab:proofs}).
\begin{table}[!htbp]
\centering
\small
\begin{tabularx}{\linewidth}{>{\raggedright\arraybackslash}p{24mm}>{\raggedright\arraybackslash}p{30mm}L}
\toprule
proof & made by, when & shows; cost \\
\midrule
spend & the sender, for each shielded transaction & its notes exist, are hers and are unspent, and the amounts balance; 6--30~s on an 8-core Mac (lichen) \\
chain & the sequencer's prover, for each sealed height & this block is a valid transition, and so was the history before it (each proof assumes the last); about 40~s for a small height, 142k cycles, on the laptop in Docker (lichen) \\
withdrawal & the user, or anyone she delegates to, only in a dispute & her statement, $P_1$ or $P_2$ (Sections~\ref{sec:withdrawal} and~\ref{sec:exit}); wrapped in Groth16 for the gadget \\
\bottomrule
\end{tabularx}
\caption{The three proofs.}
\label{tab:proofs}
\end{table}

\section{What the merged operator can do}
\label{sec:threats}

With the hub and the sequencer one party, the threat is the operator
against its users. Users against users is the ledger's ordinary job, and
the operator against itself is out of scope. Table~\ref{tab:threats}
lists what the operator could do with an unchecked account, and what
answers each.
\begin{table}[!htbp]
\centering
\small
\begin{tabularx}{\linewidth}{>{\raggedright\arraybackslash}p{58mm}L}
\toprule
the operator & answered by \\
\midrule
mints notes nobody paid for, and withdraws them from users' contracts & validity: the chain proof (Section~\ref{sec:validity}) \\
signs two valid histories that split, paying Bob in one and returning the same money to the hub in the other & one anchored history: the seal chain (Section~\ref{sec:finality}) \\
anchors a history on a branch of Bitcoin it mined & the contract's work $W$ on the anchor (Section~\ref{sec:finality}) \\
censors Alice's return or stops sealing, so that she has nothing to prove, and takes her contract when she can't & the exit (Section~\ref{sec:exit}) \\
withholds the data of something it anchored & not yet answered: data availability (Section~\ref{sec:da}) \\
\bottomrule
\end{tabularx}
\caption{What a merged operator could do, and what answers it.}
\label{tab:threats}
\end{table}
A vanished operator needs no new mechanism: a hub that never responds
to a \tx{request} lets \tx{request\_timeout} pay Alice after its window,
whatever she asserted~\cite{lichenbabe}. The exit is for one that is
present but hostile: it responds, then censors her return or stops
producing the chain receipts she needs, so that she can't prove by the
deadline and the contract is taken from her.

\section{Validity}
\label{sec:validity}

The chain guest proves one sealed height: the block's transactions
applied to the previous state give the new one, its spends authorised (a
signature, or a verified spend proof for a shielded note), no nullifier
twice, values balanced, mints signed by the hub. It takes the previous
height's chain receipt as an assumption, so each receipt vouches for the
whole history from genesis. Its journal carries the height, the state
root, a window of recent roots, the anchoring state of
Section~\ref{sec:finality}, and a history hash
\[
  \mathit{hist}_0 = 0, \qquad
  \mathit{hist}_h = H(\texttt{"lichen/hist"} \,\|\, \mathit{hist}_{h-1} \,\|\, \mathit{SR}_h),
\]
which commits the whole sequence of roots to $h$, so that an earlier
state's place in a later one's history is a hash-forward check.

Costs (lichen): a block of one mint is 395,543 cycles, a block holding
only an anchor record about 130,000, both one segment; on the laptop a
small height proves in about 40~s. The prover runs behind the seals, so
for an honest sequencer the lag is only a delay, but every withdrawal
waits for the chain receipt at its return's height. One machine proving
height after height keeps up with on the order of a hundred
transactions an hour; beyond that the remedy is the operator's: a GPU,
or proving blocks in parallel and chaining them afterwards.

A valid history is not a unique one. The sequencer can sign two
histories that split at a height, each valid, which the chain proof
cannot see. The wallet and the hub record every root they are shown and
refuse to go on when shown a second root at a height, with the evidence;
that makes equivocation visible to whoever looks, and anchoring makes it
useless.

\section{Finality}
\label{sec:finality}

\paragraph{The seal chain.} The deployment's genesis names an outpoint,
the first \emph{seal}, a P2TR output of the sequencer's key. An anchor
spends the current seal and creates the next, beside an output
\tx{OP\_RETURN "lichen"} $\| \; h \;\|\; \mathit{hist}_h$, committing the
history to a sealed height $h$; it is 269~vB (lichen). Bitcoin lets an
output be spent once, so there is one chain of anchors, after Todd's
single-use seals~\cite{toddseals}. The sequencer anchors at most once a
Bitcoin block, and only when something has changed.

\paragraph{Anchors inside the history.} A linear chain of seals could
still commit roots of two histories in turn. So once an anchor confirms,
the sequencer records it on the L2, and the chain guest checks the
record: the transaction is in the block it names, by Merkle path to the
header; it spends the seal the journal names, whose script is checked so
that no one else could spend it and stall the anchoring; and it commits
a $\mathit{hist}$ that is an ancestor of the current one. A history
forked from the anchored one can then hold no anchor after the fork: the
only spend of the seal commits the other history.

\paragraph{The contract's work.} The chain guest checks no proof of
work: an anchor on a branch of Bitcoin costs the work on top of it,
whatever the guest checked. The withdrawal statement checks it instead:
headers that follow the anchor's block, each linked and within its own
target, whose work sums to at least $W$. $W$ is fixed when the contract
opens, ten blocks' work at the difficulty of that moment (its
checkpoint, which both parties check against their own view of the
chain), and it goes into the contract's identifier,
\[
  \mathcal{D} = H(\texttt{"lichen/contract"} \,\|\, \text{deployment} \,\|\, \text{genesis seal} \,\|\, W),
\]
the deployment being the sequencer's and the hub's keys. A forger gains
nothing from headers declaring low difficulty: work is summed, so
headers worth $W$ cost about $W$ hashes however it is split. About ten
headers cost the statement a few thousand cycles. This is not a proof of
the heaviest chain; it makes a forgery a matter of mining a fork of
Bitcoin as long as $W$, which bounds the hub's exposure to one forged
history by the reward of the blocks a fork replaces.

\paragraph{For receivers.} A payment is final once its height is
anchored ten blocks deep. The wallet does not take the sequencer's word
for the anchors: it computes $\mathit{hist}$ from the roots it has
checked and finds each anchor on Bitcoin itself, so its balance splits
into final, anchored and sealed only.

\section{The withdrawal statement}
\label{sec:withdrawal}

At $P_1$ a withdrawal proves:
\begin{enumerate}[leftmargin=1.5em]
\item a note of $b$, owned by the hub with memo $c$, marked a return, in
  the commitment tree of a state root $\mathit{SR}$, by Merkle path;
\item $\mathit{SR}$ is the sequencer's, by its signature at its height;
\item a chain receipt, at that height or any later one, whose anchoring
  state commits $\mathit{SR}$: validity from genesis, in the one anchored
  history;
\item headers on the anchor's block with work at least $W$.
\end{enumerate}
Under BABE the statement also checks the hub's one-time signature on
$b$, and its journal is $c \,\|\, \mathcal{D} \,\|\, H(pk_H)$, the same for every
amount~\cite{lichenbabe}. The chain guest's image id is built into the
withdrawal guest, so the one image id the contract pins covers both.
The statement is 389,020 cycles, 482k with the hub's signature, one
segment either way (lichen); it proves in 74~s on the laptop and in
640--660~s on the server, and wraps in about 4 minutes.

\paragraph{The cooperative path.} The hub checks the same statement
natively, in milliseconds, against the chain receipt and its own node's
headers, and folds: one transaction of 193~vB. So it pays out only what
a dispute could win, and Alice proves nothing. Anchoring adds a wait,
not a cost: the next anchor and $W$ on it, about two hours on mainnet.

\section{The exit}
\label{sec:exit}

A withdrawal starts with a return, which the sequencer must include. A
merged operator profits from not including it. The exit lets Alice leave
without one:
\begin{quote}
\emph{these notes of mine were unspent at the anchored state $S$; pay me
$b$.}
\end{quote}
The operator may answer once, only with a \emph{counter}: something only
Alice could have done, showing she no longer holds what she claims. It
can't stop an exit otherwise: a censoring operator has no counter to a
true claim, and a vanished one times out.

\subsection{The claim}

Alice claims from her \emph{exit kit}, which the sequencer serves after
each of her spends and receipts and her wallet checks against the signed
root: at an anchored state $S$, each note's path in the commitment tree
and its nullifier's absence from the nullifier tree. A kit never goes
stale: the tree only grows, and absence from $S$'s set stays true of
$S$. If the operator stops serving, her last kit still works: if she
hasn't spent since, no counter exists. A payment is final for its
receiver once the receiver holds its kit.

The claim's fingerprint is
\[
  E = H(c \,\|\, O_D \,\|\, b \,\|\, h_S \,\|\, \mathit{SR}_S \,\|\, \mathit{nullifiers}),
\]
$O_D$ being the contract's deposit coin, which only its funding can
spend, so that it names the contract on Bitcoin and is known when the
setup fixes the statement. Each claimed note's owner signs $E$, and
\tx{request} posts the claim: $h_S$, $\mathit{SR}_S$, each note's opening
and its owner's signature, at most 1,169 bytes for eight notes, in
fifteen witness items of up to 80 bytes, since Bitcoin Core relays no
larger tapscript witness item. The statements read it through the
coinbase's witness commitment. It is posted so that the hub knows which
nullifiers to search for, and so that the signatures can be cited
against her. A claim spends whole notes and the contract pays once, so
the wallet claims everything it holds, up to $\vmax$; value above $b$
in the claimed notes is forfeit, and the wallet sizes its notes to fit.

\subsection{The counter}

A counter shows one of the claimed notes gone, on one of two grounds.

\paragraph{A spend.} The note's nullifier is in the nullifier tree of an
anchored state $S'$ after $S$, in the same history ($S'$'s chain receipt
hashes forward through $S$), with a leaf that counts:
\begin{itemize}[leftmargin=1.5em]
\item \emph{spent} counts if $S'$'s anchor confirmed in or before
  \tx{request}'s block. A later spend could be a transaction of hers that
  the operator held and sealed after she committed to her claim; the
  cutoff leaves only what was final on Bitcoin when she claimed.
  \tx{request} must have the contract's work on it too, or either side
  could forge the cutoff with a cheap header;
\item \emph{returned}$(c', r)$ counts the same way, unless $c' = c$ and
  $r \ge b$: then the hub holds at least $b$ of her tokens for this
  contract, which pays once. A smaller return, or one to another
  contract, can't shelter a spend;
\item \emph{exited}$(c')$ counts whenever anchored, unless $c' = c$: a
  record of the same notes claimed through another contract.
\end{itemize}

\paragraph{An equivocation.} The owner of one of the notes signed a
claim of it through another contract. Under BABE the hub responds before
Alice proves, so two requests through two contracts in one block would
both fall due before either exit record existed; her signature on the
other claim is the evidence instead. Only the owner can sign one, a
copied signature names its own contract, and an honest wallet never
signs a claim of a note through two contracts. No L2 state is involved.

\paragraph{Why the operator can't forge one.} A nullifier enters the set
only by a spend, which needs the note's keys, or by an \emph{exit
record}, which the hub posts once a \tx{request} confirms, to retire the
claimed notes, and which the chain guest accepts only with the owners'
signatures on that claim. So every counter traces to something the
owner did. Her own return to this contract never counts, so a return
she signed, which the operator still holds, can't be sealed after her
claim and used against her.

\subsection{On the gadget}

The claim and the counter ride on the dispute's first two transactions;
the rest of the graph is the main paper's~\cite{lichenbabe}.
\begin{itemize}[leftmargin=1.5em]
\item \tx{request} signs, beside $b$, the claim's $E$ under a second
  one-time key of Alice's, and posts the claim. A return's dispute signs
  a fixed $E$ and posts nothing, so one graph serves both.
\item \tx{respond} has two leaves, pre-signed at opening, both comparing
  the hub's $b$ and $E$ with Alice's digit by digit: \emph{plain}, the
  hub signs $E$ under a key $e_{\mathrm{plain}}$; \emph{countering}, it
  signs $E$ under $e_{\mathrm{counter}}$ and the counter $K$ itself.
\end{itemize}
$K$ is a Groth16 proof of the counter guest, whose journal is
$\mathcal{D} \,\|\, E$, compressed to 128 bytes and signed whole, in 259
digits. It is signed rather than hashed because no script can tie a
hash of posted bytes to signed digits without \tx{OP\_CAT}: a hub that
signed $H(K)$ could post other bytes and leave Alice nothing to prove
false. Uncompressed, at 256 bytes, it would exceed the 1,000-element
stack limit at two elements a digit.

\paragraph{The statement's three branches.} Alice's statement is fixed
at setup, with the journal $c \,\|\, \mathcal{D} \,\|\, H(pk_H \,\|\, O_D)$,
$pk_H$ the hub's four one-time keys. As $b$ does, the hub's $E$, its
mode and $K$ enter as private inputs, each authenticated by the hub's
one-time signature under a key the statement fixes, and Script checks
the same signatures in \tx{respond}. Alice wins on any of:
\begin{itemize}[leftmargin=1.5em]
\item \emph{return}: the hub signed $b$, and a return of $b$ to $c$ is
  anchored ($P_1$), in either mode;
\item \emph{claim}: the hub signed $b$ and $E$ in the plain leaf, and her
  claim holds: her notes in $S$'s tree and unspent there, signed by
  their owners, distinct, $S$ anchored with $W$ on it, \tx{request} in a
  block with $W$ on it too and posting exactly this claim, $b$ at most
  their sum;
\item \emph{failed counter}: the hub signed $E$ and $K$ in the
  countering leaf, and $K$ is not a valid counter for her
  $\mathcal{D} \,\|\, E$.
\end{itemize}
Which key signed $E$ says which leaf ran. A true counter leaves her no
branch, and \tx{respond\_\allowbreak timeout} takes the contract once the deadline
$T_A$ passes. A countered claim loses the whole contract, as any false
claim does; a careful wallet never makes one. A failed counter pays
Alice whatever her own claim: a counter is always the hub's choice, and
one that fails is one it chose to forge.

\paragraph{The hub's side.} When a claim's \tx{request} confirms, the
hub reads the claim and searches: another of its contracts' requests
signed by one of the owners, another contract's exit record, or a spend
anchored by \tx{request}'s block, its own return of at least $b$
excluded. With a ground, it assembles the counter guest's input, waits
until the blocks it needs are deep enough, proves and wraps the counter,
and responds countering. With none, it asks the sequencer to refuse
spends of the claimed nullifiers, records the exit on the L2 and
responds plainly. An honest hub may fold an exit instead, having
checked it natively.

\paragraph{What it guarantees.} Not that every redemption consumes
exactly one L2 entitlement whatever happens: a dead operator posts no
exit records, so a note can be redeemed twice. The invariant is instead
that no user loses to another or to the operator, since an honest claim
can't be countered and every payment received stays good; that every
payout not matched by value retired on the L2 is the operator's loss;
and that an operator that stays live and follows the protocol can avoid
that loss. The last fails within the anchor's latency until sender
acknowledgment (Section~\ref{sec:da}): a dishonest Alice who pays Bob
and requests before the payment's anchor confirms is paid, and Bob keeps
the payment, at the hub's cost, once per contract.

\section{What the exit costs}
\label{sec:costs}

\paragraph{The proofs.} Table~\ref{tab:cycles} gives each statement's
cycles, run on made-up inputs with the chain journal assumed (lichen),
and its proving time.
\begin{table}[!htbp]
\centering
\small
\begin{tabularx}{\linewidth}{>{\raggedright\arraybackslash}p{46mm}>{\raggedright\arraybackslash}p{24mm}L}
\toprule
statement & cycles, segments & proving \\
\midrule
a return's dispute, with the exit's guest & 489k, 1 & 640~s and a 257~s wrap on the server (lichen) \\
a claim of one note & 1.74M, 2 & 122~s on the laptop, lifted and joined (lichen) \\
a claim of eight notes & 5.27M, 6 & about 6~min on the laptop (est.) \\
the hub's counter, a spend & 594k, 1 & about a minute on the laptop (est.) \\
the hub's counter, an equivocation & 309k, 1 & 30~s on the laptop; 306~s and a 252~s wrap on the server (lichen) \\
a failed counter & 56.3M, 58 & about an hour on the laptop, and a 134~s wrap; about ten hours on the server (est., from measured segments) \\
\bottomrule
\end{tabularx}
\caption{The exit's statements: cycles, and proving times.}
\label{tab:cycles}
\end{table}
Given a valid counter, Alice's statement refuses her, both natively and
in the guest (lichen). The hub's proof is a small part of its window.
The failed counter is Alice's worst case, the subject of
Section~\ref{sec:worst}.

\paragraph{On chain.} On regtest under mainnet's relay policy (lichen):
\tx{request} 2.37~kvB with the largest claim, against 0.45~kvB without
the exit; \tx{respond} 4.48~kvB plain and 12.26~kvB countering (766
witness elements), against 0.76~kvB. So a refused true claim is about
15~kvB rather than 9.7, and an exit from a vanished hub, \tx{request}
and its timeout, about 2.9~kvB rather than 1.2 (est., from these sizes).
None of it grows with misbehaviour.

\paragraph{Windows.} The hub must find and prove a counter within its
response window $w_H$: the contract's work on \tx{request}'s block
first, then the proof and the wrap. So $w_H$ is 240 blocks on Mutinynet
(two hours), 36 on signet (six hours) and 72 on testnet4. Alice's
proof has no window: \tx{assert} is due by an absolute deadline $T_A$,
shortly before $\Tc$.

\paragraph{Congestion.} The main paper's rule, $N s \le w \cdot
1\,\mathrm{MvB}$ for every window, applies to the hub's moves. Its
windowed sum grows from 21.5 to about 25.2~kvB, or 33~kvB if every
response counters, so a dispute of a week holds about 40,000 contracts
(30,500) rather than 47,000 (est.)~\cite{lichenbabe}. Alice's windowed
move is unchanged. A wave of stale claims would make the hub prove many
counters in one $w_H$; it is costly to mount, since each countered
claimant loses a contract, and counters prove in parallel.

\section{Alice's worst case}
\label{sec:worst}

\paragraph{Why it is heavy.} The failed-counter branch must show that a
Groth16 proof does \emph{not} verify. RISC Zero composes receipts
cheaply, as the withdrawal does the chain receipt, but composition only
shows that a receipt does verify; showing that one fails means running
the verifier. $K$ has to be a Groth16 proof, not a RISC Zero receipt,
because a \tx{respond} leaf can bind only about 128 bytes of signed
digits. So the guest verifies BN254 Groth16 in software, with RISC
Zero's 256-bit accelerator: 56.3M cycles with RISC Zero's fork of
substrate-bn, 174M with arkworks (lichen). A profile of the 56.3M
(lichen): the pairings 44\%, of which Miller loops 10.7M, the final
exponentiation 8.2M and precomputation 6.6M; checking that the verifying
key's constant G2 points are in their subgroup 31\%; the proof's own
points' decompression and subgroup checks about 5.5M; the public input's
scalar multiplications 8\%.

\paragraph{What it takes.} The branch's segments prove in about 53~s
each on the laptop and take 11~s more to lift and join, so about an hour,
then 134~s to wrap. On the 4-vCPU server it would be about ten hours.
RISC Zero proves on the GPU with CUDA only; at the throughput it reports
for an RTX~4090, about 1.4~MHz~\cite{risc0bench} (not checked by us), the
segments take about a minute (est.).

\paragraph{Why it is tolerable.} Three things soften it.
\begin{itemize}[leftmargin=1.5em]
\item \emph{It is reached only by a forgery.} A hub that counters falsely
  loses the contract, so a rational hub never makes Alice run it.
\item \emph{It has days.} \tx{assert}'s deadline is absolute, close to
  the contract's end; an hour, or ten, is no constraint.
\item \emph{Anyone can prove it.} Its inputs are on chain ($E$, $K$ and
  the hub's signatures) or in the journal ($c$ and $\mathcal{D}$), so it
  needs no secret of Alice's: a proving service, her tower or a friend
  can make it, and she signs the Groth16 proof she is handed after
  checking it.
\end{itemize}
But a deterrent counts only if it is credible, and an ordinary user on a
phone cannot prove for an hour. In practice the branch presumes proving
infrastructure that someone runs for her. Every branch already asks
her for a proof and a wrap, minutes on a laptop; this one asks thirty
times more, and that is where the design is weakest.

\paragraph{Alternatives.} What we know of, and where each stands:
\begin{itemize}[leftmargin=1.5em]
\item \emph{A cheaper verifier.} Skipping the subgroup checks of the
  verifying key's constants and precomputing its constant terms would
  bring the branch to about 31M cycles (est.): half an hour on the
  laptop. A new guest, so a new deployment.
\item \emph{The equivocation ground without a SNARK.} That ground needs
  only the owner's signature, so the hub could sign the evidence itself
  instead of a proof of it. It does not touch the worst case, which the
  hub can always choose, but it shrinks the countering \tx{respond} and
  spares the hub a proof.
\item \emph{The hub signs the counter's facts, not a proof} (which note,
  how it was spent, at which height and root), and Alice refutes them
  from her own data, cheaply. Rejected: the operator can anchor a state
  it never proves, with her nullifier in it, and cite it. Against a state
  with no valid receipt Alice has nothing positive to show, whereas $K$
  proves that a valid one exists. It would also turn every withheld
  later state into a weapon against idle users, whom the design now
  protects.
\item \emph{A second BABE instance, mirrored.} The hub proves its counter
  in an instance of its own, with Alice as the garbler, and $E$, signed
  by her in \tx{request}, authenticates its input as the hub's signature
  authenticates $b$ in hers. Alice's statement loses the failed-counter
  branch; in a dispute she checks $K$ natively, in milliseconds, and
  challenges a false one. The cost moves to every contract: Alice
  garbles at each opening and renewal (11~s of computation on the
  laptop, as the hub's setup is now) and sends the hub about 240~MiB,
  and a contested counter costs her a \tx{challenge\_assert} of about
  20.7~kvB, as the hub pays now (est., from the measured setup and
  sizes). With contracts of months that is a modest price, and the one
  alternative that removes Alice's proving altogether.
\end{itemize}

\section{Open questions}
\label{sec:open}

\subsection{Data availability}
\label{sec:da}

The exit is safe exactly when Alice can read the anchored data she
needs. Anchors are public; their contents are not. The sharpest form is
the \emph{consolidation attack}: Alice consolidates her notes into one,
as her wallet recommends; the operator anchors the consolidation and
withholds the new note's data. Her old notes are countered, the
consolidation really spent them, and the new one has no kit. The same
works on any spend that leaves her change. If nothing has been anchored
since her kit, she is safe.

A counter wins the hub a contract, so withholding is harmless whenever a
spend whose data Alice can't read can't be cited. Most remedies are
rules of that kind~\cite{bitcoinrollups}:
\begin{itemize}[leftmargin=1.5em]
\item \emph{Sender acknowledgment}, the planned one. A spend can be cited
  only with its sender's signature that she holds its data, which her
  wallet gives once it has her kit at the anchored state containing the
  spend; until then the spend's outputs are locked, and only she may
  claim them. No trust and nothing on Bitcoin, at the price of
  two-phase finality: Bob's payment is final once the acknowledgment is
  anchored too. It also closes the race of Section~\ref{sec:exit}, since
  the cutoff can then allow a grace the other way.
\item \emph{Publishing each anchor's inserted nullifiers and new
  commitments}, about 24~vB per one-in-two-out transaction: no trust,
  a cost per transaction on Bitcoin.
\item \emph{A data committee in the seal}: the seal's key a threshold key
  of the operator and a committee, each member signing an anchor only
  once it holds the batch's data. One honest signer suffices; the price
  is the committee's liveness for progress.
\item \emph{Accept the gap}: a user is exit-safe only while the operator
  serves data for her latest transactions.
\end{itemize}

\subsection{The rest}

\begin{itemize}[leftmargin=1.5em]
\item The exit played on chain with real proofs, and the failed-counter
  branch proved whole, not only estimated from other inputs' segments.
\item Alice's worst case: whether the mirrored instance, or something
  better, should replace the failed-counter branch
  (Section~\ref{sec:worst}).
\item Shielded notes' exits. A shielded note's authority is knowledge of
  $\mathit{nk}$ inside the proof, with nothing to post or cite, so today
  a wallet exits its transparent notes only. Shielded addresses would
  gain a BIP340 owner key, bound into the nullifier, whose signature on
  $E$ is posted like a transparent owner's; an exit would then reveal
  that key and link the wallet's exits, though not its spends.
\item A tower for the exit. Alice's proof depends on which leaf the hub
  responds with, so unlike a return's she cannot sign it ahead; she
  must come back after \tx{respond}, or entrust a key.
\item A proof of the heaviest chain, or a counter-proof of a heavier one,
  in place of the work threshold $W$.
\item Mass exit: every contract claiming by $\Tc$ against a hostile
  operator. The counter adds no games, since a failed counter pays
  Alice, but larger requests and a longer schedule, which favour few
  notes per claim, and so consolidation, which needs data availability.
\end{itemize}

\begin{thebibliography}{99}

\bibitem{lichen} waxwing. Lichen: L2s in channels enforced by Bitcoin. Version 0.3, October 2026. \url{https://github.com/AdamISZ/lichen/blob/master/papers/lichen-v0.3.1.pdf}

\bibitem{lichenbabe} waxwing. Lichen over BABE: withdrawals of a variable amount, by authenticated request. Version 0.1, October 2026. \url{https://github.com/AdamISZ/lichen/blob/master/papers/lichen-babe-v0.1.9.pdf}

\bibitem{lichencode} lichen: the ledger, the sequencer, the hub and the wallet. \url{https://github.com/AdamISZ/lichen}

\bibitem{shielded} lichen's shielded transfers: an example L2. \url{https://github.com/AdamISZ/lichen/blob/master/docs/SHIELDED.md}

\bibitem{babe} S.~Garg, D.~Kolonelos, M.~Sergeevitch, S.~Sridhar and D.~Tse. BABE: verifying proofs on Bitcoin made 1000x cheaper. ePrint 2026/065, revised October 2026. \url{https://eprint.iacr.org/2026/065}

\bibitem{bitvm2usenix} R.~Linus Woll, L.~Aumayr, Z.~Avarikioti, M.~Maffei, A.~Pelosi, O.~S.~Thyfronitis Litos, C.~Stefo, D.~Tse and A.~Zamyatin. Bridging Bitcoin to second layers via BitVM2. USENIX Security 2026; ePrint 2025/1158. \url{https://eprint.iacr.org/2025/1158}

\bibitem{toddseals} P.~Todd. Preventing consensus fraud with commitments and single-use-seals, 2016. \url{https://petertodd.org/2016/commitments-and-single-use-seals}

\bibitem{bitcoinrollups} J.~Light. Validity rollups on Bitcoin, Appendices C and D, 2022. \url{https://bitcoinrollups.org}

\bibitem{risc0bench} RISC Zero. zkVM benchmarks and datasheet. \url{https://dev.risczero.com/api/zkvm/benchmarks}

\end{thebibliography}

\end{document}
